\section{Idea 5: Checkpoint-Fork Querying}
\label{sec:checkpoint-fork}

\subsection{Thesis and novelty boundary}

An uncertain label is not necessarily useful: either answer may induce nearly
the same policy.  \forkquery\ asks for feedback only when the two possible
labels produce virtual one-step checkpoints with different behavior on
user-relevant probe states.

\Inherited: \prlhf\ and \pdpo, active preference collection, AdamW, LoRA,
Jensen--Shannon divergence, and information-gain acquisition.

\Proposed: score a query in \emph{behavior space} after both exact
label-conditioned optimizer branches, then commit only the observed branch.
The finite probe score is not a guarantee of downstream utility.

\begin{table}[t]
  \centering
  \caption{Algorithm papers supporting Checkpoint-Fork Querying.}
  \label{tab:fork-support}
  \begin{tabularx}{\textwidth}{@{}L{0.24\textwidth}Y Y@{}}
    \toprule
    Supporting paper & Component used here & Missing component supplied by this proposal \\
    \midrule
    P-RLHF \citep{li2024personalized} &
    Personalized DPO loss and shared/private parameterization. &
    It trains a fixed preference set and never allocates live queries. \\
    SELM \citep{zhang2024selm} &
    Active exploration for online language-model alignment. &
    It does not value both labels by the behavior of their resulting
    checkpoints. \\
    Hybrid Preference Optimization \citep{bose2024hybrid} &
    Joint use of offline preference data and online queries. &
    It does not define a decision-focused per-query acquisition score. \\
    SPRInG \citep{kim2026spring} &
    Selective updates under continual personalization. &
    Likelihood novelty does not test whether two labels change future actions. \\
    Information-gain acquisition \citep{mackay1992information} &
    Select data by expected reduction of decision uncertainty. &
    Classical information gain does not include a functional LLM optimizer
    branch or agent action probes. \\
    \bottomrule
  \end{tabularx}
\end{table}

\subsection{Step-by-step workflow}

\begin{enumerate}
  \item Train and load an immutable offline \prlhf\ parent.
  \item In each user window, generate a fixed pool of same-task live trajectory
  pairs with the exact active checkpoint.
  \item Keep only pairs matched on deterministic task success and safety, so
  the requested label targets personalization rather than obvious correctness.
  \item Before scoring the pool, freeze one anchor minibatch, user-relevant
  probe states, and schema-valid action slates.
  \item For each candidate and each possible label, clone the parent adapter
  and AdamW state functionally and execute the exact one-step \pdpo\ update.
  \item Evaluate both virtual checkpoints on every fixed action slate.
  \item Compute the label-prior-weighted Jensen--Shannon divergence between
  their action distributions and choose the largest fork.
  \item Abstain if the maximum score is below the calibrated threshold.
  Otherwise request one label and execute only its branch as a real update from
  the unchanged parent.
  \item Hash, load, validate, activate, or roll back the real candidate.  All
  virtual branches and unqueried labels expire.
\end{enumerate}

\subsection{Virtual branches and fork score}

For candidate pair $q=(u,x,y^A,y^B)$, define
\begin{equation}
  d_{\params,e}(q)
  =
  \log\frac{\pi_{\params}(y^A\mid x,e)}
            {\refpolicy(y^A\mid x,e)}
  -
  \log\frac{\pi_{\params}(y^B\mid x,e)}
            {\refpolicy(y^B\mid x,e)}.
\end{equation}
For hypothetical label $z\in\{+1,-1\}$, where $+1$ prefers $y^A$, use
\begin{equation}
  \ell_z(q;\params)
  =
  -\log\sigma\!\left(\beta z d_{\params,e_u}(q)\right)
  -
  \lambda_g\log\sigma\!\left(\beta z d_{\params,e_0}(q)\right).
  \label{eq:fork-pdpo}
\end{equation}
Freeze a seven-pair eligible anchor minibatch $H_t$ for the whole acquisition
window and set
\begin{equation}
  \loss_{q,z}
  =
  \frac{1}{8}\left[
    \ell_z(q;\params)+\sum_{h\in H_t}\ell(h;\params)
  \right].
\end{equation}
From the same parent parameters and AdamW state $(m,v,r)$, each hypothetical
branch computes
\begin{align}
  g_{q,z} &= \nabla_{\params}\loss_{q,z}, \\
  m'_{q,z} &= \beta_1m+(1-\beta_1)g_{q,z}, \\
  v'_{q,z} &= \beta_2v+(1-\beta_2)g_{q,z}^2, \\
  \widetilde\params_{q,z}
  &=
  \params-\alpha\left[
    \frac{\widehat m'_{q,z}}{\sqrt{\widehat v'_{q,z}}+\varepsilon}
    +\lambda_{\mathrm{wd}}\params
  \right].
  \label{eq:fork-branch}
\end{align}

For each of $R=8$ training-side probe states $s_r$, fix a slate
$A_r=\{a_{r1},\ldots,a_{rK}\}$ of $K=4$ valid actions.  The branch
distribution is
\begin{equation}
  P_{q,z,r}(j)
  =
  \softmax_j\!\left(
    |a_{rj}|^{-1}
    \log\pi_{\widetilde\params_{q,z}}(a_{rj}\mid s_r,e_u)
  \right).
\end{equation}
Calibrate the label prior $p_q=\sigma(d_{\params,e_u}(q)/T_c)$ only on
eligible historical labels, and let
$M_{q,r}=p_qP_{q,+,r}+(1-p_q)P_{q,-,r}$.  The acquisition score is
\begin{equation}
  F(q)
  =
  \frac{1}{R}\sum_{r=1}^{R}
  \left[
    p_q\KL(P_{q,+,r}\Vert M_{q,r})
    +(1-p_q)\KL(P_{q,-,r}\Vert M_{q,r})
  \right].
  \label{eq:fork-score}
\end{equation}
This is the mutual information between the binary label and the finite action
slate under the approximation above.

\begin{algorithm}[t]
  \caption{Checkpoint-Fork Querying}
  \label{alg:checkpoint-fork}
  \KwIn{Parent $c_0$; candidate-pool size $M$; anchors $H_t$; probes and
  slates; threshold $\delta$; validation panel $V$}
  \KwOut{Queried labels, abstentions, and activated checkpoints}
  $c\leftarrow c_0$\;
  \While{label budget, compute budget, and patience remain}{
    collect $M$ same-task live pairs with active checkpoint $c$\;
    filter to matched task-success and safety status\;
    freeze $H_t$, probes, slates, and calibrated label priors\;
    \ForEach{candidate $q$}{
      \ForEach{$z\in\{-1,+1\}$}{
        clone the same parent adapter and AdamW state functionally\;
        compute virtual branch $\widetilde\params_{q,z}$ by \cref{eq:fork-branch}\;
        evaluate fixed action slates and discard the virtual state\;
      }
      compute $F(q)$ by \cref{eq:fork-score}\;
    }
    $q^\star\leftarrow\arg\max_qF(q)$\;
    \eIf{$F(q^\star)<\delta$}{
      record an abstention; expire the pool\;
    }{
      request label $z^\star$ only for $q^\star$\;
      from the unchanged parent, take one real AdamW step on $\loss_{q^\star,z^\star}$\;
      save and hash candidate $\hat c$\;
      \eIf{$\hat c$ loads and passes $V$}{
        activate $\hat c$; $c\leftarrow\hat c$\;
      }{
        restore $c$ and its optimizer state\;
      }
    }
  }
\end{algorithm}

\subsection{Invariants, stopping, and tests}

\begin{itemize}
  \item Both virtual labels start from the same parent parameters, moments,
  anchors, probes, learning rate, and weight decay.
  \item Virtual branches leave persistent parameters and moments bitwise
  unchanged and can never be activated.
  \item Labels remain hidden until a candidate has been selected and its query
  receipt written.
  \item Probe states come from eligible training-side history, never validation
  or untouched future tasks.
  \item Every acquisition baseline sees the same pool, hidden labels, feedback
  budget, and real-update implementation.
\end{itemize}

A user window stops after one query and real update, an abstention, or its
rollout/token ceiling.  Required acquisition baselines are random selection,
label entropy, SPRInG-style novelty, expected gradient norm, expected model
change, and a compute-matched sham that executes both branches and probes but
selects randomly.  Parameter-space fork distance and first-order logit
approximations test whether full behavior-space simulation is necessary.

\Evidence.  The method may spend substantial compute to avoid a single weak
label, and its score depends on probe coverage and prior calibration.  The
primary empirical question is personalization per acquired label and per
activated update, with task success, safety, and total GPU cost reported
separately.
